\documentclass[11pt,a4paper]{article}
\usepackage[margin=1in]{geometry}
\usepackage[T1]{fontenc}
\usepackage[utf8]{inputenc}
\usepackage{lmodern}
\usepackage{microtype}
\usepackage{amsmath,amssymb,amsthm,mathtools,bm}
\usepackage{booktabs,array,longtable}
\usepackage[hidelinks]{hyperref}
\usepackage{enumitem}
\usepackage{physics}
\usepackage{setspace}
\setstretch{1.06}

% SST macro prelude
\newcommand{\swirlarrow}{\mathchoice{\mkern-2mu\scriptstyle\boldsymbol{\circlearrowleft}}{\mkern-2mu\scriptstyle\boldsymbol{\circlearrowleft}}{\mkern-2mu\scriptscriptstyle\boldsymbol{\circlearrowleft}}{\mkern-2mu\scriptscriptstyle\boldsymbol{\circlearrowleft}}}
\newcommand{\swirlarrowcw}{\mathchoice{\mkern-2mu\scriptstyle\boldsymbol{\circlearrowright}}{\mkern-2mu\scriptstyle\boldsymbol{\circlearrowright}}{\mkern-2mu\scriptscriptstyle\boldsymbol{\circlearrowright}}{\mkern-2mu\scriptscriptstyle\boldsymbol{\circlearrowright}}}
\newcommand{\vswirl}{\mathbf{v}_{\swirlarrow}}
\newcommand{\vswirlcw}{\mathbf{v}_{\swirlarrowcw}}
\newcommand{\SwirlClock}{S_t^{\swirlarrow}}
\newcommand{\SwirlClockcw}{S_t^{\swirlarrowcw}}
\newcommand{\rhoF}{\rho_{\!f}}
\newcommand{\rhoE}{\rho_{\!E}}
\newcommand{\rhoM}{\rho_{\!m}}
\newcommand{\rhocore}{\rho_{\text{core}}}
\newcommand{\rc}{r_c}

\title{Toward a Canonical Atomic Torsion-Bridge in Swirl-String Theory\\
\large v0.7 Circulation-Foundational Atomic Torsion Bridge}
\author{Omar Iskandarani}
\date{March 2026}

\begin{document}
\maketitle

\begin{abstract}
This note consolidates the minimal reusable core of the recent atomic torsion-bridge drafts into a circulation-foundational v0.7 using the retained SST canons. The central change is ontological and mathematical: the primitive circulation quantum
\[
\Gamma_0 = 2\pi r_c\,\lVert \mathbf{v}_{\!\boldsymbol{\circlearrowleft}}\rVert
\]
is promoted above the clock shift, so that the atomic sector is driven first by local circulation and only secondarily by the Swirl-Clock branches. The retained ingredients are: (i) a propagating transverse torsion/shear field represented by a divergence-free vector potential, (ii) a rotating-fluid source law that converts axial displacement gradients into vertical vorticity and azimuthal swirl, (iii) a canon-level circulation normalization through $\Gamma_0$, and (iv) a branch-resolved atomic Hamiltonian in which $\SwirlClock$ and $\SwirlClockcw$ remain the internal two-valued branch. The conservative Coulombic atomic limit and Hartree screening are retained. The new SST content is concentrated in a circulation-sensitive interaction consisting of a scalar branch shift plus a circulation-orbit coupling proportional to an anticommutator with $\hat L_z$. A photoionization benchmark is formulated by coupling the torsion field to continuum states through a branch-resolved Fermi-golden-rule transition rate, and the frozen interaction yields an operator-level handedness asymmetry already for the hydrogenic $1s$ source state. The intentionally excluded pieces remain excluded: the ultra-fast compressional branch, skyrmion-emission analogies, and magneto-optical extensions are left outside the canonical core until independently secured.
\end{abstract}

\section{Aim and canonical constraints}
The objective is not to replace the successful mathematics of atomic physics, but to reinterpret it through a minimal SST field theory while preserving its tested limits. The required observable layer is therefore retained as
\begin{equation}
\label{eq:state-labels}
\text{electron state}=(n,\ell,m_\ell,\sigma),
\qquad
\sigma\in\{\swirlarrow,\swirlarrowcw\},
\end{equation}
with the intended low-energy identification that the SST branch label $\sigma$ reproduces the observed two-valued internal sector ordinarily encoded by $m_s=\pm \tfrac12$.

A minimally canonical SST atomic equation must satisfy four constraints:
\begin{enumerate}[label=(C\arabic*)]
    \item recover the hydrogenic spectrum in the one-electron weak-coupling limit;
    \item preserve the shell-state counting $N_{n\ell}=2(2\ell+1)$ and $N_n=2n^2$;
    \item admit a closed many-electron mean-field interpretation via a self-consistent density $\rho$;
    \item isolate new SST content to explicit, testable couplings rather than hidden redefinitions.
\end{enumerate}

\section{Retained ingredients from the earlier drafts}
Three ingredients are retained.

\paragraph{(i) Propagating torsion projection.} The earlier photon/torsion draft identifies an electromagnetic-strength two-form as a projection of torsion,
\begin{equation}
\label{eq:Ftorsion}
F \equiv \kappa\,u_a T^a,
\end{equation}
in a teleparallel setting, and models the propagating radiation sector as a transverse director/shear wave with
\begin{equation}
\label{eq:uA}
\mathbf{u}=\partial_t \mathbf{A},
\qquad
\nabla\cdot \mathbf{A}=0,
\end{equation}
so that, with $\mathbf u$ a velocity field, $\mathbf A$ carries dimensions of length. The retained wave equation is the corresponding massless transverse field equation.

\paragraph{(ii) Rotating-fluid source law.} The rotating-cylinder drafts provide a concrete source chain in a uniformly rotating incompressible fluid. For axisymmetric perturbations,
\begin{equation}
\label{eq:wz}
\partial_t \omega_z = 2\Omega\,\partial_z w,
\qquad
w=\partial_t \xi,
\end{equation}
which integrates from rest to
\begin{equation}
\label{eq:wzxi}
\omega_z(r,z,t)=2\Omega\,\partial_z \xi(r,z,t).
\end{equation}
The azimuthal swirl then follows kinematically from
\begin{equation}
\label{eq:utheta}
\omega_z = \frac{1}{r}\partial_r(r u_\theta)
\quad\Longrightarrow\quad
u_\theta(r,z,t)=\frac{1}{r}\int_0^r \omega_z(r',z,t)\,r'\,dr'.
\end{equation}
This yields a reversible leading-order swirl response for symmetric up-down forcing.

\paragraph{(iii) Swirl-Clock bookkeeping.} The SST Rosetta draft cleanly distinguishes the material scale
\begin{equation}
\label{eq:Omega0}
\Omega_0 \equiv \frac{\lVert \vswirl\rVert}{\rc},
\end{equation}
from the laboratory rotation rate $\Omega$, and retains the clock factor
\begin{equation}
\label{eq:SwirlClock}
\SwirlClock \equiv S_t^{(\swirlarrow)} = \sqrt{1-\frac{v_{\swirlarrow}^{2}}{c^2}},
\qquad
\SwirlClockcw \equiv S_t^{(\swirlarrowcw)} = \sqrt{1-\frac{v_{\swirlarrowcw}^{2}}{c^2}},
\end{equation}
as the canonical time-scaling variable. In the branch-symmetric limit one sets $v_{\swirlarrow}=v_{\swirlarrowcw}=v$, so that the two clocks coincide; away from that limit the arrow labels track the two internal branches of the atomic sector.

\paragraph{(iv) Primitive circulation quantum.} The retained canons elevate quantized circulation to a primitive SST ingredient,
\begin{equation}
\label{eq:Gamma0}
\boxed{
\Gamma_0 \equiv 2\pi \rc\,\lVert \vswirl\rVert,
\qquad
\Gamma = \oint \mathbf v\cdot d\boldsymbol\ell = n\,\Gamma_0.
}
\end{equation}
This means the natural canon-first hierarchy is circulation $\to$ clock, not clock $\to$ circulation. In the same canons, delay-selected circulation modes precede knot selection, and Kelvin protection of circulation is treated as foundational rather than optional \cite{IskandaraniCanon078,IskandaraniCanon060}. The present v0.7 draft therefore uses $\Gamma_0$ as the primary normalization of the atomic driver rather than the earlier purely scalar clock shift.

\section{Minimal retained field sector}
The first retained sector is a transverse torsion/shear field. The minimal Lagrangian is
\begin{equation}
\label{eq:Ltorsion}
\mathcal{L}_{\rm torsion}
=
\frac12\rho_{\rm eff}\,\lVert \partial_t \mathbf{A}\rVert^2
-
\frac12 K\,\lVert \nabla\times\mathbf{A}\rVert^2,
\qquad
\nabla\cdot\mathbf{A}=0.
\end{equation}
Varying with respect to $\mathbf{A}$ gives
\begin{equation}
\label{eq:Awave-pre}
\rho_{\rm eff}\left(\partial_t^2\mathbf{A}+c^2\nabla\times\nabla\times\mathbf{A}\right)=0,
\qquad
c^2\equiv \frac{K}{\rho_{\rm eff}},
\end{equation}
which reduces, under $\nabla\cdot\mathbf{A}=0$, to
\begin{equation}
\label{eq:Awave}
\boxed{
\left(\frac{1}{c^2}\partial_t^2-\nabla^2\right)\mathbf{A}=0.
}
\end{equation}
This is the propagating field retained from the older torsion/photon sector. No magneto-optical or parity-odd vacuum extensions are retained in the present canonical core.

\section{Minimal retained source sector}
The second retained sector is the rotating-medium source law. The atomic theory proposed below does not require that every torsion excitation originate in a macroscopic rotating cylinder, but the rotating-cylinder calculation provides the cleanest explicit production mechanism for an SST circulation amplitude.

The circulation-foundational choice retained in v0.7 is to normalize the source by local circulation, not by local clock shift. Using Eq.~\eqref{eq:utheta}, define the local loop circulation by
\begin{equation}
\label{eq:GammaLoc}
\boxed{
\Gamma_{\rm loc}(r,z,t)
:=
2\pi r\,u_\theta(r,z,t)
=
2\pi\int_0^r \omega_z(r',z,t)\,r'\,dr'.
}
\end{equation}
The canon-level dimensionless circulation driver is then
\begin{equation}
\label{eq:gammaDef}
\boxed{
\gamma(\mathbf r,t)
\equiv
\frac{\Gamma_{\rm loc}(\mathbf r,t)}{\Gamma_0}.
}
\end{equation}
Equivalently, via Eq.~\eqref{eq:wzxi},
\begin{equation}
\label{eq:gammaXi}
\gamma(r,z,t)
=
\frac{4\pi\Omega}{\Gamma_0}
\int_0^r \partial_z \xi(r',z,t)\,r'\,dr'.
\end{equation}
Thus the new source chain is
\begin{equation}
\label{eq:source-chain}
\xi \longrightarrow \omega_z \longrightarrow u_\theta \longrightarrow \Gamma_{\rm loc} \longrightarrow \gamma,
\end{equation}
with the clock sector treated downstream of circulation rather than as the primitive driver.

For a localized helical packet, a practical canon-compatible ansatz is
\begin{equation}
\label{eq:gammaPacket}
\gamma(r,\phi,z,t)
=
A_\gamma
\exp\!\left(-\frac{r^2}{2w_r^2}\right)
\exp\!\left(-\frac{(z-vt)^2}{2w_z^2}\right)
\cos\!\left(m\phi+kz-\omega t\right).
\end{equation}
This profile is retained because it carries finite amplitude, axial propagation, helical winding, and finite support, while remaining directly interpretable as a circulation-bearing packet.

For continuity with earlier versions, the scalar vorticity-normalized driver
\begin{equation}
\label{eq:ThetaDef}
\Theta(\mathbf r,t) \equiv \frac{\omega_z(\mathbf r,t)}{\Omega_0}
\end{equation}
may still be used as an auxiliary bookkeeping field. In the present v0.7 draft, however, it is no longer the primary canonical source variable.

\section{Atomic sector: branch-resolved minimal Hamiltonian}
The atomic bridge model is built by retaining the ordinary Coulomb core and many-electron screening, and placing the novel SST content entirely in a circulation-sensitive branch potential. The working Hamiltonian is
\begin{equation}
\label{eq:master-H}
\boxed{
\left[
-\frac{\hbar^2}{2\mu}\nabla^2
+V_{\rm core}(r)
+V_{\rm screen}[\rho](\mathbf r)
+V_{\rm circ}^{(\sigma)}[\gamma]
\right]\Psi_{i,\sigma}
=\varepsilon_{i,\sigma}\Psi_{i,\sigma},
}
\end{equation}
with
\begin{equation}
\label{eq:Vcore}
V_{\rm core}(r)=-\frac{Ze^2}{4\pi\varepsilon_0 r}
\end{equation}
for the conservative baseline.

A provisional ``freeze v0.5'' removes the free branch couplings in favor of canonical SST scales. First, define the branch-coupling strength by the small dimensionless material ratio
\begin{equation}
\label{eq:eta0}
\boxed{
\eta_0 \equiv \frac{\lVert\vswirl\rVert}{c} = \frac{\alpha}{2} \approx 3.648676\times 10^{-3}.
}
\end{equation}
Then take the unperturbed atomic baseline to be branch-symmetric,
\begin{equation}
\label{eq:St0freeze}
S_{t,0}(r)\equiv 1,
\end{equation}
and define the two branch sectors by a circulation-driven clock bookkeeping field
\begin{equation}
\label{eq:StBranches}
\boxed{
S_t^{(\sigma)}(\mathbf r,t)=1+\sigma\,\eta_0\,\gamma(\mathbf r,t),
\qquad
\sigma=\pm 1,
}
\end{equation}
where $\sigma=+1$ corresponds to $\SwirlClock$ and $\sigma=-1$ to $\SwirlClockcw$. The clock is therefore retained, but only as a downstream representation of circulation-bearing dynamics.

Next, freeze the energy scale to the Rydberg binding scale,
\begin{equation}
\label{eq:ERfreeze}
\boxed{
E_R \equiv \frac{\hbar^2}{2\mu (a_0^{\rm SST})^2}
= \frac{\mu e^4}{2(4\pi\varepsilon_0)^2\hbar^2}.
}
\end{equation}
In v0.4 the clock potential also carried a gradient piece proportional to $(a_0^{\rm SST})^2\abs{\nabla S_t^{(\sigma)}}^2$. In the present v0.7 freeze that term is dropped deliberately rather than silently: once circulation is promoted to the primitive atomic driver, the spatial structure previously delegated to the gradient penalty is instead carried by the explicit circulation field $\gamma(\mathbf r,t)$ and by the orbit-coupled anticommutator term introduced below. The frozen starter model therefore keeps only the algebraic clock shift in $V_{\rm clock}^{(\sigma)}$ and treats spatial/chiral structure as belonging to the circulation sector.

The scalar clock-derived contribution is then
\begin{equation}
\label{eq:Vclock}
\boxed{
V_{\rm clock}^{(\sigma)}
=
E_R\Bigl[1-\bigl(S_t^{(\sigma)}\bigr)^2\Bigr].
}
\end{equation}
At first order this gives
\begin{equation}
\label{eq:Vclock-linear}
\boxed{
\delta V_{\rm clock}^{(\sigma)}
\approx
-2\sigma E_R\eta_0\,\gamma.
}
\end{equation}

The genuinely circulation-foundational addition retained in v0.7 is an orbit-coupled term. To avoid introducing a new fitted constant, the same small canonical factor $\eta_0$ is reused and the coupling is frozen at the same atomic energy scale. The minimal Hermitian choice is the anticommutator form
\begin{equation}
\label{eq:HintCirc}
\boxed{
\delta H_{\rm circ}^{(\sigma)}
=
\frac{E_R\eta_0\,\sigma}{2}
\left\{\frac{\hat L_z}{\hbar},\gamma\right\}.
}
\end{equation}
This form is algebraically equivalent to the previous $\{\hat L_z,\gamma\}/(2\hbar)$ representation, but makes the scaling explicit: $\hat L_z/\hbar$ is dimensionless, so the anticommutator is dimensionless before multiplication by the common frozen atomic energy scale $E_R\eta_0$. The anticommutator is important: for an initial $s$-state with $\hat L_z\psi_{1s}=0$, the operator still yields a nonzero contribution because $\hat L_z$ can act on the helical packet factor inside $\gamma$.

The full frozen linear interaction is therefore
\begin{equation}
\label{eq:HintTotal}
\boxed{
\delta H_{\rm int}^{(\sigma)}
=
-2\sigma E_R\eta_0\,\gamma
+
\frac{E_R\eta_0\,\sigma}{2}
\left\{\frac{\hat L_z}{\hbar},\gamma\right\}.
}
\end{equation}
This is the main circulation-foundational upgrade relative to the earlier scalar-clock baseline: the atomic sector now responds directly to handed circulation, not only to a scalar branch shift.

\section{Many-electron closure and screening}
To reach shell structure, a many-electron mean-field closure is required. The electron density is defined by
\begin{equation}
\label{eq:rho-def}
\rho(\mathbf r)=\sum_{i,\sigma} f_{i,\sigma}\,\abs{\Psi_{i,\sigma}(\mathbf r)}^2,
\end{equation}
with occupation numbers $f_{i,\sigma}\in\{0,1\}$ at the effective level. The minimal retained screening functional is the Hartree term,
\begin{equation}
\label{eq:Vscreen}
\boxed{
V_{\rm screen}[\rho](\mathbf r)
=
\frac{e^2}{4\pi\varepsilon_0}
\int \frac{\rho(\mathbf r')}{\abs{\mathbf r-\mathbf r'}}\,d^3r'.
}
\end{equation}
No SST-specific correction to screening is introduced in the canonical starter model. This is deliberate: a predictive theory must first recover the ordinary many-electron limit before attempting branch- or medium-induced corrections to screening.

At this stage Pauli antisymmetry is retained as an effective low-energy rule rather than claimed as derived. The shell-state counts then remain
\begin{equation}
\label{eq:subshell-count}
N_{n\ell}=2(2\ell+1),
\end{equation}
and
\begin{equation}
\label{eq:shell-count}
N_n=\sum_{\ell=0}^{n-1}2(2\ell+1)=2n^2.
\end{equation}
In the present bridge model, the factor $2\ell+1$ arises from the ordinary rotational sectors $m_\ell=-\ell,\ldots,+\ell$, while the factor $2$ is interpreted as the low-energy branch doublet associated with $\SwirlClock$ and $\SwirlClockcw$.

\section{Hydrogenic limit and SST length scale}
The model must recover the ordinary hydrogenic spectrum when
\begin{equation}
V_{\rm screen}\to 0,
\qquad
\gamma\to 0,
\qquad
S_t^{(+)}=S_t^{(-)}.
\end{equation}
Then~\eqref{eq:master-H} reduces to the standard one-electron Coulomb problem,
\begin{equation}
\label{eq:En}
E_n=-\frac{\mu Z^2 e^4}{2(4\pi\varepsilon_0)^2\hbar^2}\frac{1}{n^2}.
\end{equation}
The SST atomic scale is tied to the canonical constants by
\begin{equation}
\label{eq:a0-sst}
\boxed{
a_0^{\rm SST}
=
\frac{F_{\rm swirl}^{\max}\,\rc^2}{M_e\,\lVert\vswirl\rVert^2}
=
\frac{2\rc}{\alpha^2}.
}
\end{equation}
Using the canonical values
\begin{equation}
\lVert\vswirl\rVert = 1.09384563\times 10^6\ \mathrm{m\,s^{-1}},
\quad
\rc = 1.40897017\times 10^{-15}\ \mathrm{m},
\quad
F_{\rm swirl}^{\max}=29.053507\ \mathrm{N},
\end{equation}
this gives
\begin{equation}
\label{eq:a0-num}
a_0^{\rm SST}
=
5.29177161\times 10^{-11}\ \mathrm{m},
\end{equation}
while the equivalent relation $2\rc/\alpha^2$ yields
\begin{equation}
\label{eq:a0-num2}
\frac{2\rc}{\alpha^2}=5.29177214\times 10^{-11}\ \mathrm{m},
\end{equation}
consistent with the ordinary Bohr scale to the expected rounding level.

This gives the hydrogenic SST radius law
\begin{equation}
\label{eq:RnZ}
R_{n,Z}^{\rm SST}\sim \frac{n^2}{Z}\,a_0^{\rm SST}.
\end{equation}
Thus the atomic bridge model retains the tested size law while reinterpreting $a_0$ as a medium-derived coherence scale. In the present draft this agreement is to be read as a \emph{consistency check on the retained canonical scales}, not as an independent prediction, since $(F_{\rm swirl}^{\max},\rc,\lVert\vswirl\rVert)$ were fixed upstream in the SST program.

\section{Transition amplitudes and torsion capture}
A circulation-bearing torsion packet~\eqref{eq:gammaPacket} drives atomic transitions through the frozen interaction~\eqref{eq:HintTotal}. For an initial state $\ket{i,\sigma}$ and final state $\ket{f,\sigma'}$, the first-order matrix element is
\begin{equation}
\label{eq:Mfi}
M_{fi}^{\sigma'\sigma}=\matrixel{f,\sigma'}{\delta H_{\rm int}}{i,\sigma}.
\end{equation}
The conservative resonance condition is
\begin{equation}
\label{eq:resonance}
\hbar\omega \approx \varepsilon_f-\varepsilon_i.
\end{equation}
If the packet carries azimuthal phase factor $e^{im\phi}$, then the orbital projection rule remains
\begin{equation}
\label{eq:dm}
\Delta m_\ell = m.
\end{equation}
The new feature is that the anticommutator term in Eq.~\eqref{eq:HintCirc} is explicitly odd under $m\mapsto -m$. For the helical packet ansatz of Eq.~\eqref{eq:gammaPacket}, the radial and axial Gaussian envelopes are $\phi$-independent, so $\hat L_z=-i\hbar\partial_\phi$ acts only on the azimuthal factor. One therefore has the exact identity
\begin{equation}
\label{eq:Lzgamma}
\hat L_z\,\gamma = m\hbar\,\gamma,
\end{equation}
so the interaction can distinguish opposite-handed packets at operator level rather than only after selecting final states. This is exactly the chiral handle absent from the scalar-only v0.4 freeze.

\paragraph{Operator-level $1s$ asymmetry remark.}
For the hydrogenic ground state, $\hat L_z\psi_{1s}=0$. Applying Eq.~\eqref{eq:HintTotal} together with Eq.~\eqref{eq:Lzgamma} gives
\begin{equation}
\label{eq:Hint1s}
\boxed{
\delta H_{\rm int}^{(\sigma)}\psi_{1s}
=
E_R\eta_0\,\sigma\left(-2+\frac{m}{2}\right)\gamma\,\psi_{1s}.
}
\end{equation}
Thus the helical packet channels are already inequivalent before continuum integration:
\begin{equation}
\label{eq:Hint1sPm}
 m=+1:\quad \delta H_{\rm int}^{(\sigma)}\psi_{1s}=-\frac{3}{2}E_R\eta_0\,\sigma\,\gamma\,\psi_{1s},
 \qquad
 m=-1:\quad \delta H_{\rm int}^{(\sigma)}\psi_{1s}=-\frac{5}{2}E_R\eta_0\,\sigma\,\gamma\,\psi_{1s}.
\end{equation}
If the corresponding continuum overlap integrals satisfy $|I_{+1}(k,\omega)|=|I_{-1}(k,\omega)|$ to leading order, then the source-level rate asymmetry is already nonzero:
\begin{equation}
\label{eq:A1sSource}
\boxed{
\mathcal A_{1s}^{\rm(source)}
=
\frac{|{-3}/{2}|^2-|{-5}/{2}|^2}{|{-3}/{2}|^2+|{-5}/{2}|^2}
=
-\frac{8}{17}
\approx -0.4706.
}
\end{equation}
Equation~\eqref{eq:A1sSource} is not yet the full continuum photoionization result; it is the operator-level asymmetry implied by the frozen interaction before the Coulomb-continuum matrix elements are evaluated.

This suggests the following interpretation. Upward shell excitation under resonant incoming torsion is the atomic analogue of a ``giant jet,'' while downward relaxation accompanied by outgoing torsion radiation is the analogue of a ``sprite.'' These are analogical labels only; the retained mathematics is ordinary time-dependent perturbation theory acting on a branch-resolved Hamiltonian.

\section{Continuum sector and photoionization benchmark}
The sharpest next benchmark is photoionization rather than further shell bookkeeping. A branch-resolved atomic theory is incomplete until it specifies not only its bound spectrum but also its continuum and the dynamical channel by which an incoming torsion packet transfers energy into that continuum. The minimal conservative choice is to define the continuum as the positive-energy sector of the same branch-resolved Hamiltonian already introduced in Eq.~\eqref{eq:master-H}.

Because $S_{t,0}\equiv 1$ in the present freeze, the static clock contribution vanishes identically. The unperturbed branch Hamiltonian is therefore simply
\begin{equation}
\label{eq:H0sigma}
\boxed{
H_0^{(\sigma)}
=
-\frac{\hbar^2}{2\mu}\nabla^2
+
V_{\rm core}(r)
+
V_{\rm screen}[\rho](\mathbf r),
\qquad
\sigma=\pm 1.
}
\end{equation}
The corresponding effective potential is
\begin{equation}
\label{eq:Veffsigma}
V_{\rm eff}^{(\sigma)}(\mathbf r)
=
V_{\rm core}(r)
+
V_{\rm screen}[\rho](\mathbf r).
\end{equation}
Bound states satisfy
\begin{equation}
\label{eq:boundstates}
H_0^{(\sigma)}\psi_{n\ell m_\ell}^{(\sigma)}
=
E_{n\ell}^{(\sigma)}\psi_{n\ell m_\ell}^{(\sigma)},
\qquad
E_{n\ell}^{(\sigma)}<0,
\end{equation}
whereas continuum scattering states satisfy
\begin{equation}
\label{eq:continuumstates}
H_0^{(\sigma)}\psi_{\mathbf k}^{(\sigma)}
=
E_{\mathbf k}^{(\sigma)}\psi_{\mathbf k}^{(\sigma)},
\qquad
E_{\mathbf k}^{(\sigma)}>0.
\end{equation}
Under the localized-packet assumption
\begin{equation}
\label{eq:gammaAsymptotic}
\gamma(\mathbf r,t)\to 0 \qquad (\lVert \mathbf r\rVert\to\infty),
\end{equation}
we have $S_t^{(\sigma)}\to 1$ and hence no asymptotic branch-dependent offset. The continuum threshold is therefore branch-degenerate and may be taken as
\begin{equation}
\label{eq:continuum-threshold}
E_{\rm cont}^{(+)}(0)=E_{\rm cont}^{(-)}(0)=0.
\end{equation}

The next required step is to close the dynamical link between the propagating torsion field and the circulation driver $\gamma$. The retained v0.7 closure is to extract vorticity from the field velocity $\mathbf u=\partial_t\mathbf A$ and then build the loop circulation from that vorticity. Define
\begin{equation}
\label{eq:omegaA}
\omega_z^{(A)}(\mathbf r,t)
:=
\hat{\mathbf z}\cdot\left(\nabla\times \partial_t \mathbf A\right)(\mathbf r,t),
\end{equation}
then
\begin{equation}
\label{eq:GammaA}
\Gamma_{\rm loc}^{(A)}(r,z,t)
:=
2\pi\int_0^r \omega_z^{(A)}(r',z,t)\,r'\,dr',
\end{equation}
and finally
\begin{equation}
\label{eq:gammaA}
\boxed{
\gamma^{(A)}(\mathbf r,t)
=
\frac{\Gamma_{\rm loc}^{(A)}(\mathbf r,t)}{\Gamma_0}.
}
\end{equation}
This is the key circulation-foundational replacement for the older scalar-insertion strategy. The field now drives the atom through canon-normalized circulation.

To permit back-reaction one must source the field equation itself. The branch-sensitive sourced wave equation is retained as
\begin{equation}
\label{eq:Awave-source}
\boxed{
\left(\frac{1}{c^2}\partial_t^2-\nabla^2\right)\mathbf A
=
g_J \,\widetilde{\mathbf J}_{\rm br},
}
\end{equation}
with the source length frozen to the canonical core scale
\begin{equation}
\label{eq:gJfreeze}
\boxed{
 g_J \equiv r_c = \frac{\alpha^2}{2}\,a_0^{\rm SST}.
}
\end{equation}
The corresponding branch current is
\begin{equation}
\label{eq:Jbr}
\boxed{
\mathbf J_{\rm br}
=
\frac{\hbar}{2 i \mu}
\sum_{\sigma=\pm 1}
\sigma
\left(
\Psi_\sigma^*\nabla\Psi_\sigma
-
(\nabla\Psi_\sigma^*)\Psi_\sigma
\right),
\qquad
\widetilde{\mathbf J}_{\rm br}\equiv \frac{\mathbf J_{\rm br}}{\Omega_0}.
}
\end{equation}
Equation~\eqref{eq:Jbr} remains a provisional closure ansatz rather than a theorem of SST, but it keeps the field-atom-source loop closed with no free prefactor.

Once $\gamma$ is defined by Eq.~\eqref{eq:gammaA}, the frozen circulation interaction from Eq.~\eqref{eq:HintTotal} becomes the genuine photoionization operator. The branch-resolved ionization threshold from an initial bound state $i$ to a final continuum branch $\sigma_f$ is
\begin{equation}
\label{eq:IonizationThresholdGeneral}
I_i^{(\sigma_i\to \sigma_f)}
=
E_{\rm cont}^{(\sigma_f)}(0)-E_i^{(\sigma_i)}.
\end{equation}
Under the asymptotic choice~\eqref{eq:continuum-threshold}, this reduces to
\begin{equation}
\label{eq:IonizationThresholdSimple}
I_i^{(\sigma_i)}=-E_i^{(\sigma_i)}.
\end{equation}
The photoionization condition is then frequency-thresholded,
\begin{equation}
\label{eq:PhotoThreshold}
\boxed{
\hbar\omega \ge I_i^{(\sigma_i\to \sigma_f)}.
}
\end{equation}
For an emitted free electron with asymptotic kinetic energy $K$, one obtains
\begin{equation}
\label{eq:KineticEnergyAtom}
\boxed{
K
=
\hbar\omega - I_i^{(\sigma_i\to \sigma_f)}.
}
\end{equation}
For a condensed system with work function $\phi^{(\sigma_i\to \sigma_f)}$, the analogous relation is
\begin{equation}
\label{eq:PhotoelectricSolid}
\boxed{
K_{\max}
=
\hbar\omega - \phi^{(\sigma_i\to \sigma_f)}.
}
\end{equation}
Thus the threshold is controlled by frequency, while the intensity enters only through the occupation number of the incoming torsion mode.

To show this explicitly, let $N_\omega$ denote the occupation number of an incoming torsion packet of frequency $\omega$. Then the continuum version of Fermi's golden rule becomes
\begin{equation}
\label{eq:GoldenRuleContinuum}
\boxed{
\Gamma_{i\to \mathbf k,\sigma_f}
=
\frac{2\pi}{\hbar}
N_\omega
\left|
\matrixel{\psi_{\mathbf k}^{(\sigma_f)}}{\delta H_{\rm int}^{(\sigma_f)}}{\psi_i^{(\sigma_i)}}
\right|^2
\delta\!\left(
E_{\mathbf k}^{(\sigma_f)}-E_i^{(\sigma_i)}-\hbar\omega
\right).
}
\end{equation}
This is the required rate structure if the SST torsion packet is to reproduce the key empirical logic of the photoelectric effect: an immediate threshold in frequency, and an emission rate proportional to flux or intensity rather than a threshold controlled by amplitude.

A useful distinction therefore emerges between two regimes. In the \emph{weak-field quantum absorption regime}, Eqs.~\eqref{eq:PhotoThreshold}--\eqref{eq:GoldenRuleContinuum} define the canonical baseline. In a possible later \emph{strong-field nonlinear regime}, sufficiently large coherent $\gamma$ could reshape the effective barrier and produce field-emission or over-the-barrier escape. That second regime should be treated as an extension and not conflated with the baseline photoelectric benchmark.

The branch-sensitive content now enters through an explicitly handed operator. Define the handedness observable
\begin{equation}
\label{eq:Asymmetry}
\boxed{
\mathcal A(\omega)
=
\frac{\Gamma_{h=+1}-\Gamma_{h=-1}}{\Gamma_{h=+1}+\Gamma_{h=-1}},
}
\end{equation}
where $h=\pm 1$ labels the handedness of the incoming torsion packet. For the helical packet ansatz~\eqref{eq:gammaPacket} at fixed propagation direction, the handedness is identified with the sign of the azimuthal winding number,
\begin{equation}
\label{eq:h-sign-m}
\boxed{h\equiv \mathrm{sgn}(m).}
\end{equation}
Combined with the orbital rule $\Delta m_\ell=m$ from Eq.~\eqref{eq:dm}, the anticommutator term in Eq.~\eqref{eq:HintCirc} implies that same-handed channels acquire an operator-level factor proportional to $m$. The corresponding starter-model overlap preference is
\begin{equation}
\label{eq:handed-selection}
 h=+1:\ \sigma_i=\sigma_f=+1\ 	ext{favored},
 \qquad
 h=-1:\ \sigma_i=\sigma_f=-1\ 	ext{favored},
\end{equation}
with opposite-handed channels suppressed by reduced helical overlap and branch-flip channels absent at linear order. In contrast to v0.4, this preference is now traceable directly to the interaction operator rather than only to packet labeling.

If $\mathcal A(\omega)=0$ identically even after including Eq.~\eqref{eq:HintCirc}, then the circulation-foundational extension still collapses to a relabeling. If instead $\mathcal A(\omega)\neq 0$ with a definite symmetry and scale, then the theory acquires a genuinely SST-specific atomic signal. This is the main computational fork opened by the present v0.7 draft.

\section{What is intentionally left outside the canonical core}
Three earlier ingredients are intentionally excluded from the present article.

First, the ultra-fast compressional branch with kinematic front speed $c_P\gg c$ is excluded. It is too loosely coupled to the retained atomic sector and too constrained by external bounds to serve as part of a first canonical core.

Second, the skyrmion-emission and topological-inheritance material is excluded from the canonical starter model. It is topologically interesting but not required to close the atomic Hamiltonian.

Third, vacuum magneto-optical extensions such as parity-odd Faraday rotation are excluded. They may become downstream consequences of the torsion field sector, but they are not needed to define the atomic bridge itself.

\section{Minimal predictive commitments}
The present model earns the label ``possibly canonical'' only if its commitments are kept narrow. The retained predictive commitments are:
\begin{enumerate}[label=(P\arabic*)]
    \item hydrogenic recovery: Eq.~\eqref{eq:En} must emerge when $\gamma\to 0$ and $V_{\rm screen}\to 0$;
    \item shell counting: the low-energy atomic sector must reproduce Eqs.~\eqref{eq:subshell-count} and~\eqref{eq:shell-count};
    \item branch symmetry: if $\gamma=0$, the two branches must be degenerate;
    \item circulation sensitivity: if $\gamma\neq 0$, then branch-sensitive line shifts or transition asymmetries are allowed and, in principle, measurable;
    \item frozen starter couplings: in the present v0.5 freeze, $\Gamma_0$, $\gamma$, $\eta_0$, $g_J$, and $\delta H_{\rm int}^{(\sigma)}$ are given explicitly by Eqs.~\eqref{eq:Gamma0}, \eqref{eq:gammaDef}, \eqref{eq:eta0}, \eqref{eq:gJfreeze}, and~\eqref{eq:HintTotal}, with no atom-by-atom retuning.
\end{enumerate}

The strongest falsifier is simple. If no branch-sensitive effects survive once the couplings are fixed by canon, then the branch interpretation reduces to a relabeling of ordinary spin without new atomic content. If branch-sensitive effects do survive, their magnitude and symmetry scaling become direct experimental targets.

\section{Status}
This article does not claim a first-principles derivation of fermionic exchange, nor a first-principles derivation of the branch-resolved clock sector beyond the provisional freeze given above. It is therefore not yet a closed atomic theory. What it does provide is a sharply reduced and internally consistent circulation-foundational starter model:
\begin{equation}
\boxed{
\text{one propagating torsion field}
+\text{ one rotating-fluid source law}
+\text{ one branch-resolved atomic Hamiltonian}.
}
\end{equation}
That reduction is its main contribution. It turns three partially overlapping drafts into one narrow framework that can be tested, extended, or discarded term by term.

\paragraph{Analogy for a 10-year-old.} Imagine the atom as a globe with allowed storm patterns on it. A twisting wave packet can hit the globe and nudge one storm pattern into a higher allowed one. The two Swirl-Clock branches are like two opposite ways the storm can twist internally. The point of this paper is to say exactly which equations describe the wave, the circulation injection, and the storm pattern.

\begin{thebibliography}{99}

\bibitem{IskandaraniSST10}
O. Iskandarani,
\textit{Impulsive Axisymmetric Forcing in a Rotating Cylinder, Reversible Swirl Response, and Skyrmionic Photon Fields: Fluid Benchmarks and Fluid-Inspired Kinematic Hypotheses},
Zenodo preprint, 2026.
DOI: 10.5281/zenodo.17619170.

\bibitem{IskandaraniSST11}
O. Iskandarani,
\textit{Impulsive Axisymmetric Forcing in a Rotating Cylinder, Reversible Swirl Response, and Skyrmionic Photon Emission: Fluid Benchmarks and Fluid-Inspired Kinematic Hypotheses},
Zenodo preprint, 2026.
DOI: 10.5281/zenodo.17619573.

\bibitem{IskandaraniSST17}
O. Iskandarani,
\textit{Electromagnetism as Propagating Torsion in a Hydrodynamic Vacuum: A Geometric Unification via Cartan Structure Equations},
Zenodo preprint, 2026.
DOI: 10.5281/zenodo.17677074.

\bibitem{IskandaraniCanon078}
O. Iskandarani,
\textit{Swirl-String Theory Canon v0.7.8},
unpublished internal technical reference, 2026.

\bibitem{IskandaraniCanon060}
O. Iskandarani,
\textit{Swirl-String Theory Canon v0.6.0},
unpublished internal technical reference, 2025.

\bibitem{Greenspan1968}
H. P. Greenspan,
\textit{The Theory of Rotating Fluids},
Cambridge University Press, 1968.

\bibitem{Batchelor1967}
G. K. Batchelor,
\textit{An Introduction to Fluid Dynamics},
Cambridge University Press, 1967.

\bibitem{Vallis2017}
G. K. Vallis,
\textit{Atmospheric and Oceanic Fluid Dynamics},
2nd ed., Cambridge University Press, 2017.
DOI: 10.1017/9781107588417.

\bibitem{GriffithsSchroeter2018}
D. J. Griffiths and D. F. Schroeter,
\textit{Introduction to Quantum Mechanics},
3rd ed., Cambridge University Press, 2018.

\bibitem{BetheSalpeter1957}
H. A. Bethe and E. E. Salpeter,
\textit{Quantum Mechanics of One- and Two-Electron Atoms},
Springer, 1957.

\bibitem{Dirac1928}
P. A. M. Dirac,
``The Quantum Theory of the Electron,''
\textit{Proceedings of the Royal Society A} \textbf{117} (1928), 610--624.
DOI: 10.1098/rspa.1928.0023.

\bibitem{HehlObukhov2007}
F. W. Hehl and Y. N. Obukhov,
\textit{Elie Cartan's Torsion in Geometry and in Field Theory, an Essay},
Annales de la Fondation Louis de Broglie \textbf{32} (2007), 157--194.
Permalink: hal.science/hal-00408133.

\end{thebibliography}

\end{document}
